% ============================================================
% pragmatic-style.tex — Stile Pragmatic Bookshelf
% ============================================================
% Caratteristiche: font moderno, margini con annotazioni,
% codice con numeri riga, colori caldi (arancio/marrone),
% capitoli con banda colorata.
% ============================================================

% === LINGUA ===
\usepackage[\BabelLang]{babel}
\usepackage[autostyle]{csquotes}

% === BIBLIOGRAFIA ===
\usepackage[
  backend=biber,
  style=numeric,
  sorting=nyt,
  urldate=long,
  maxbibnames=3,
  backref=true,
  backrefstyle=none
]{biblatex}

% === FONT (Pragmatic: Bookman/Century, monospace stretto) ===
\usepackage{fontspec}
\usepackage{amsmath}
\usepackage{amssymb}
\setmainfont{texgyrebonum}[
  Extension=.otf,
  UprightFont=*-regular,
  BoldFont=*-bold,
  ItalicFont=*-italic,
  BoldItalicFont=*-bolditalic
]
\setsansfont{texgyreadventor}[
  Extension=.otf,
  Scale=0.90,
  UprightFont=*-regular,
  BoldFont=*-bold,
  ItalicFont=*-italic,
  BoldItalicFont=*-bolditalic
]
\setmonofont{InconsolataN}[
  Extension=.otf,
  Scale=0.85,
  UprightFont=*-Regular,
  BoldFont=*-Bold
]

% === LAYOUT PAGINA (7.5 × 9 poll. — formato Pragmatic) ===
\usepackage[
  paperwidth=7.5in,
  paperheight=9in,
  top=0.85in, bottom=0.85in,
  inner=0.95in, outer=1.1in,
  marginparwidth=0.8in,
  marginparsep=0.15in,
  headheight=14pt,
  footskip=0.45in
]{geometry}

% === SPAZIATURA PARAGRAFI ===
\setlength{\parindent}{1.2em}
\setlength{\parskip}{0pt plus 1pt}

% === COLORI (palette Pragmatic — arancio/terracotta) ===
\usepackage[dvipsnames,svgnames,table]{xcolor}
\definecolor{prag-orange}{HTML}{D35400}
\definecolor{prag-dark}{HTML}{2C3E50}
\definecolor{prag-brown}{HTML}{6D4C41}
\definecolor{prag-gray}{HTML}{7F8C8D}
\definecolor{prag-lightgray}{HTML}{BDC3C7}
\definecolor{code-bg}{HTML}{FFFBE6}
\definecolor{code-border}{HTML}{E0D8B0}
\definecolor{tip-green}{HTML}{27AE60}
\definecolor{warn-orange}{HTML}{E67E22}
\definecolor{note-blue}{HTML}{2980B9}
\definecolor{link-blue}{HTML}{1A5276}
\definecolor{danger-red}{HTML}{C0392B}

% === GRAFICA E TABELLE ===
\usepackage{graphicx}
\usepackage{float}
\usepackage{booktabs}
\usepackage{longtable}
\usepackage{tabularx}
\usepackage{multirow}
\usepackage{array}
\usepackage{colortbl}

% === CODICE SORGENTE (Pragmatic: sfondo caldo, numeri riga) ===
\usepackage{listings}
\lstset{
  basicstyle=\footnotesize\ttfamily\color{prag-dark},
  backgroundcolor=\color{code-bg},
  frame=single,
  framerule=0.5pt,
  rulecolor=\color{code-border},
  breaklines=true,
  breakatwhitespace=true,
  showstringspaces=false,
  tabsize=2,
  xleftmargin=2em,
  xrightmargin=0.5em,
  aboveskip=0.8em,
  belowskip=0.8em,
  numbers=left,
  numberstyle=\tiny\color{prag-lightgray},
  numbersep=8pt,
  keepspaces=true,
  keywordstyle=\color{prag-orange}\bfseries,
  stringstyle=\color{tip-green},
  commentstyle=\color{prag-gray}\itshape,
  columns=fullflexible,
  literate={€}{{\euro}}1 {→}{{$\rightarrow$}}1 {←}{{$\leftarrow$}}1
}
\lstdefinelanguage{yaml}{
  keywords={true,false,null,yes,no},
  sensitive=false,
  comment=[l]{\#},
  morestring=[b]",
  morestring=[b]',
}

% === INTESTAZIONI (Pragmatic: capitolo e sezione, minimale) ===
\usepackage{fancyhdr}
\pagestyle{fancy}
\fancyhf{}
\fancyhead[LE]{\small\color{prag-gray}\leftmark}
\fancyhead[RO]{\small\color{prag-gray}\rightmark}
\fancyfoot[LE,RO]{\small\sffamily\color{prag-brown}\thepage}
\renewcommand{\headrulewidth}{0pt}

% === TITOLI CAPITOLO (Pragmatic: banda arancio + numero) ===
\usepackage{titlesec}
\titleformat{\chapter}[display]
  {\normalfont\sffamily\bfseries}
  {\colorbox{prag-orange}{\parbox[c][1.2cm][c]{1.5cm}{\centering\fontsize{36}{36}\selectfont\color{white}\thechapter}}}
  {14pt}
  {\huge\color{prag-dark}}
\titlespacing*{\chapter}{0pt}{-10pt}{25pt}
\titleformat{\section}
  {\normalfont\Large\sffamily\bfseries\color{prag-dark}}
  {\color{prag-orange}\thesection}{0.7em}{}
\titleformat{\subsection}
  {\normalfont\large\sffamily\bfseries\color{prag-brown}}
  {\thesubsection}{0.7em}{}
\titleformat{\subsubsection}
  {\normalfont\normalsize\sffamily\bfseries\color{prag-gray}}
  {\thesubsubsection}{0.7em}{}

% === TITOLI PARTI ===
\titleformat{\part}[display]
  {\centering\normalfont\sffamily\bfseries}
  {\fontsize{48}{48}\selectfont\color{prag-orange!30}\partname\ \thepart}
  {12pt}
  {\Huge\color{prag-dark}}

% === HYPERLINK ===
\usepackage{xurl}
\usepackage[
  colorlinks=true,
  linkcolor=prag-dark,
  citecolor=prag-orange,
  urlcolor=link-blue,
  bookmarks=true,
  bookmarksnumbered=true
]{hyperref}

% === PROTEZIONE OVERFLOW ===
\tolerance=1000
\emergencystretch=1.5em
\setlength{\tabcolsep}{4pt}

% === SPAZIATURA LISTE ===
\usepackage{enumitem}
\setlist{nosep, leftmargin=1.8em, topsep=0.3em, itemsep=0.15em}

% === MARGINNOTES (Pragmatic signature feature) ===
\usepackage{marginnote}
\renewcommand*{\marginfont}{\footnotesize\sffamily\color{prag-gray}}

% === BOX ADMONITION ===
% ============================================================
% pragmatic-boxes.tex — Box admonition stile Pragmatic Bookshelf
% ============================================================
% Stile: box con bordo colorato, angoli arrotondati, icone.
% Pragmatic usa: Tip, Note, Important, Joe Asks, Did You Know?
% ============================================================

\providecolor{tip-green}{HTML}{27AE60}
\providecolor{warn-orange}{HTML}{E67E22}
\providecolor{note-blue}{HTML}{2980B9}
\providecolor{danger-red}{HTML}{C0392B}
\providecolor{prag-brown}{HTML}{6D4C41}
\providecolor{prag-gray}{HTML}{7F8C8D}

\usepackage{fontawesome5}
\usepackage{tcolorbox}
\tcbuselibrary{skins,breakable}

% Tip
\newtcolorbox{tipbox}[1][]{
  colback=tip-green!5, colframe=tip-green,
  fonttitle=\sffamily\bfseries,
  title={\faLightbulb\space \LabelTip},
  boxrule=1pt, arc=4pt, breakable,
  left=6pt, right=6pt,
  coltitle=white, colbacktitle=tip-green, #1
}

% Note
\newtcolorbox{notebox}[1][]{
  colback=note-blue!5, colframe=note-blue,
  fonttitle=\sffamily\bfseries,
  title={\faCommentDots\space \LabelNote},
  boxrule=1pt, arc=4pt, breakable,
  left=6pt, right=6pt,
  coltitle=white, colbacktitle=note-blue, #1
}

% Warning / Important
\newtcolorbox{warnbox}[1][]{
  colback=warn-orange!5, colframe=warn-orange,
  fonttitle=\sffamily\bfseries,
  title={\faExclamationCircle\space \LabelImportant},
  boxrule=1pt, arc=4pt, breakable,
  left=6pt, right=6pt,
  coltitle=white, colbacktitle=warn-orange, #1
}

% Danger
\newtcolorbox{dangerbox}[1][]{
  colback=danger-red!5, colframe=danger-red,
  fonttitle=\sffamily\bfseries,
  title={\faTimesCircle\space \LabelDanger},
  boxrule=1pt, arc=4pt, breakable,
  left=6pt, right=6pt,
  coltitle=white, colbacktitle=danger-red, #1
}

% "Joe Asks" style — question/curiosity box
\newtcolorbox{checkpointbox}[1][]{
  colback=prag-brown!5, colframe=prag-brown,
  fonttitle=\sffamily\bfseries,
  title={\faQuestionCircle\space \LabelDidYouKnow},
  boxrule=1pt, arc=4pt, breakable,
  left=6pt, right=6pt,
  coltitle=white, colbacktitle=prag-brown, #1
}

% Version
\newtcolorbox{versionbox}[1][]{
  colback=purple!4, colframe=purple!60!black,
  fonttitle=\sffamily\bfseries,
  title={\faCogs\space \LabelVersion},
  boxrule=1pt, arc=4pt, breakable,
  left=6pt, right=6pt,
  coltitle=white, colbacktitle=purple!60!black, #1
}

% Cost
\newtcolorbox{costbox}[1][]{
  colback=Goldenrod!5, colframe=Goldenrod!80!black,
  fonttitle=\sffamily\bfseries,
  title={\faCoins\space \LabelCost},
  boxrule=1pt, arc=4pt, breakable,
  left=6pt, right=6pt,
  coltitle=white, colbacktitle=Goldenrod!80!black, #1
}

% Exercise
\newtcolorbox{praticabox}[1][]{
  colback=prag-gray!5, colframe=prag-gray,
  fonttitle=\sffamily\bfseries,
  title={\faPencilAlt\space \LabelExercise},
  boxrule=1pt, arc=4pt, breakable,
  left=6pt, right=6pt,
  coltitle=white, colbacktitle=prag-gray, #1
}


% === EURO ===
\newcommand{\euro}{€}

% === INDICE ANALITICO ===
\usepackage{makeidx}
\makeindex
